% Part: applied-modal-logics
% Chapter: temporal-logic
% Section: introduction

\documentclass[../../../include/open-logic-section]{subfiles}

\begin{document}

\olfileid{aml}{tl}{int}

\olsection{પરિચય}

સમયસંબંધી તર્કપદ્ધતિઓ એવી બાબતો વિશેના દાવાઓનું અધ્યયન કરે છે જે
આગળ સાચી થશે અથવા અગાઉ સાચી હતી. સમયસંબંધી તર્કના પ્રણેતા તરીકે
આર્થર પ્રાયરને શ્રેય આપવામાં આવે છે; તેમણે તેને \emph{કાળ-તર્ક}
કહ્યો હતો. અહીં અમે સમયસંબંધી તર્કનો અભ્યાસ મોટાભાગે પ્રાયરના મૂળ
મોડલ અભિગમ મુજબ કરીશું: તેમાં વિધાનતર્કના મૂળભૂત માળખામાં
સમયસંબંધી સંક્રિયકો ઉમેરાય છે, જ્યારે વિધાનતર્ક સામાન્ય રીતે
વિધાનોને કાળ-નિર્દેશ વિનાનાં ગણે છે.

દાખલા તરીકે, વિધાનતર્કમાં હું બીઝી નામના કૂતરા વિશે વાત કરી શકું:
કૂતરાંઓની ટેવ મુજબ તે ક્યારેક બેસે છે અને ક્યારેક બેસતો નથી.
શાસ્ત્રીય તર્કમાં બીઝી બેસે છે અને બીઝી બેસતો નથી, એવા બંને દાવા
એકસાથે કરવા વિરોધાભાસી ગણાય. પરંતુ બંને દાવા સાચા હોઈ શકે છે,
માત્ર એક જ સમયે નહીં; ભાષામાં સમયસંબંધી સંક્રિયકો ઉમેરવાથી આ વાત
સરળતાથી વ્યક્ત કરી શકાય છે. આ સંક્રિયકો ઉમેરવાથી ``બીઝીને આગળ
ખાવાનું ઇનામ અથવા દડો મળશે'' પરથી ``બીઝીને આગળ ખાવાનું ઇનામ
મળશે અથવા બીઝીને આગળ દડો મળશે'' જેવા અનુમાનોની માન્યતા પણ
સમજાવી શકાય છે.

જોકે સમયસંબંધી તર્કમાં ઘણા દાર્શનિક પ્રશ્નો ઊભા થાય છે, જેને કારણે
આપણે તેની એક રૂપરેખા બીજી કરતાં પસંદ કરી શકીએ. દાખલા તરીકે,
ભવિષ્યસંબંધી આકસ્મિક વિધાન ભવિષ્ય વિશેનું એવું વિધાન છે જે ન તો
અનિવાર્ય છે ન તો અશક્ય. જો આપણે કહીએ કે ``રિચાર્ડ આવતી કાલે
કરિયાણાની દુકાને જશે'', તો આપણે હજી ન બનેલી બાબત વિશે દાવો કરીએ
છીએ, જેનું સત્યમૂલ્ય વિવાદાસ્પદ છે. હકીકતમાં, એ દાવાને મૂળે
સત્યમૂલ્ય \emph{આપી} પણ શકાય છે કે કેમ એ વિવાદાસ્પદ છે. જો આપણે
કડક નિયતિવાદી હોઈએ, તો સંબંધિત ઘટના બનવાની હોય તે પહેલાં પણ આ
વાક્ય હકીકતમાં સાચું કે ખોટું છે, એવું માનવામાં કદાચ આપણને વાંધો
ન હોય---માત્ર તેનું સત્યમૂલ્ય હજી આપણે જાણતા ન હોઈએ. તેનાથી
ઊલટું, આપણે એવું માની શકીએ કે ભવિષ્ય ખરેખર ખુલ્લું છે અને
ભવિષ્યસંબંધી આકસ્મિક વિધાનોના સત્યમૂલ્યો અનિર્ધારિત છે.

સમયની રચના અને સ્વરૂપ અંગેની આવી ઘણી માન્યતાઓ સમયસંબંધી
તર્કપદ્ધતિઓનાં નિદર્શનો અને રૂપરેખાઓની આપણી પસંદગીમાં વણાયેલી હોય છે.
દાખલા તરીકે, આપણે પૂછવું પડે કે સમય રેખીય, શાખાવાળો કે ચક્રીય
હોય એવા નિદર્શનો રચવા જોઈએ? આપણાં સમયસંબંધી નિદર્શનોમાં આરંભબિંદુ અને
અંતબિંદુ હોવા જોઈએ કે નહીં, અને સમયને વિવિક્ત ક્ષણો દ્વારા
દર્શાવવો કે સાતત્ય તરીકે, એ પણ નક્કી કરવું પડે.

\end{document}
